\documentclass[11pt,a4paper]{article}
\usepackage[a4paper,margin=25mm]{geometry}
\usepackage[UTF8]{ctex}
\usepackage{microtype}
\usepackage[hidelinks]{hyperref}
\setlength{\parindent}{0pt}
\setlength{\parskip}{0.75em}
\begin{document}

\begin{flushright}
夏昂然\\
陕西师范大学化学化工学院\\
陕西省西安市长安区西长安街620号 710119\\
\href{mailto:haangyeon@snnu.edu.cn}{haangyeon@snnu.edu.cn}\\[1em]
\today
\end{flushright}

尊敬的\emph{Genes}期刊编辑：

您好！

谨此提交我们的研究论文"跨模型不确定性感知的顺式调控元件最小编辑：细胞类型选择性设计"，申请作为研究论文发表。本研究聚焦于功能基因组学中的一个核心生物学问题：如何用少量、可解释的核苷酸替换重调天然顺式调控元件（CRE），同时保留细胞类型选择性的调控信号。

现有的序列设计工作流已使序列-活性预测达到很高的效率，但通常使用定义目标的同一个预测器来选择编辑，或者返回一条无约束的高分序列。SafeEdit-CRE提出了一套程序分离的最小编辑框架：候选生成使用一个预测器，跨模型评审使用架构差异的集成模型，最终评估使用留出的多核模型。该框架将已发表的细胞类型特异性预测器与验证校准的不确定性、序列域控制以及天然200-nt CRE骨架上的约束束搜索相结合，优先选择在独立训练、架构不同的模型间预测一致的变异体。

本研究的主要发现包括：
\begin{itemize}
  \item 在K562、HepG2和SK-N-SH三种细胞的600个留出天然CRE亲本上，SafeEdit-CRE将跨模型特异性裕度增益从贪婪编辑的0.822提升至0.877（配对差值0.055；95\% CI 0.040--0.071）。
  \item 与计算量匹配的主模型束搜索相比，SafeEdit-CRE保持了高度吻合的预测特异性（差值$-0.004$；95\% CI $-0.011$--$0.004$），同时跨模型不确定性从0.335降至0.314。
  \item 在锁定的非重叠90亲本设计质量基准中，通过全部预设序列和模型一致性检查的比例从贪婪编辑的38.8\%提升至SafeEdit-CRE的60.6\%，产出74个优先用于报告基因筛选的Tier A变异体。
\end{itemize}

这些结果为\emph{Genes}的期刊范围提供了聚焦的贡献：它们将序列建模与顺式调控语法、细胞类型特异性转录调控以及可重复的功能基因组学设计联系起来。完整代码、候选文库和冻结汇总输入已整理为v1.1.0发布包，将在正式发表前同步至GitHub和Zenodo；同行评审期间可向通信作者索取，最终DOI将在发表版本中补入。

本文为原创研究，未在其他期刊考虑发表，已获得作者批准。作者声明无利益冲突。感谢您对本研究的考虑。

此致\\[0.8em]
夏昂然

\end{document}
